\documentclass[10pt,a4paper]{amsart}
\usepackage[margin=19mm]{geometry}
\usepackage{amsmath,amssymb,mathtools}
\usepackage[scr=rsfs,cal=euler]{mathalfa}
\usepackage{xfrac}
\usepackage[unicode,hidelinks,bookmarks=false]{hyperref}
\newcommand{\ie}{i.e.\ }
\newcommand{\cf}{cf.\ }
\newcommand{\resp}{respectively\ }
\newcommand{\Subsection}{\subsection}
\providecommand{\nobreakdash}{\nobreak\hbox{-}}
\setlength{\emergencystretch}{2em}
\multlinegap0pt
\usepackage{bm}
\usepackage{bbm}
\usepackage{dsfont}
\usepackage{xspace}
\usepackage{cancel}
\usepackage{mathtools}
\usepackage{cleveref}
\usepackage[shortlabels]{enumitem}
\usepackage{amsthm}
\makeatletter
\@ifundefined{theorem}{\newtheorem{theorem}{Theorem}}{}
\@ifundefined{lemma}{\newtheorem{lemma}[theorem]{Lemma}}{}
\makeatother
\DeclareMathOperator{\SL}{SL}
\DeclareMathOperator{\rank}{rk} 
\newcommand{\CC}{\mathbb{C}} 
\newcommand{\I}{I}
\newcommand{\J}{J}
\newcommand{\K}{K}
\newcommand{\PP}{\mathbb{P}} 
\newcommand{\cA}{\mathcal{A}}
\newcommand{\cB}{\mathcal{B}} 
\newcommand{\overbar}[1]{\mkern 1.5mu\overline{\mkern-1.5mu#1\mkern-1.5mu}\mkern 1.5mu}
\renewcommand{\L}{L}
\renewcommand{\subset}{\subseteq}
\title[Tschirnhausen Bundles of Quintic Covers of $\mathbb{P}^1$]{Tschirnhausen Bundles of Quintic Covers of $\mathbb{P}^1$}
\author{Sam Frengley, Sameera Vemulapalli}
\date{Source: arXiv:2507.06942v4 (CC BY 4.0)}
\begin{document}
\maketitle

\noindent\emph{Source and licence.} Tschirnhausen Bundles of Quintic Covers of $\mathbb{P}^1$. Version arXiv:2507.06942v4 (CC BY 4.0), distributed under the Creative Commons Attribution 4.0 International licence, as recorded in the arXiv licence metadata for this exact version and shown on the abstract page https://arxiv.org/abs/2507.06942v4. Full rights notice: licenses/arxiv-2507.06942-CC-BY-4.0.txt.\par\smallskip

\noindent\textit{Editorial scope.} Counted unit: the Lemma whose source label is \texttt{lemma:divisibility} in the article (source0.tex lines 520--527, source label \texttt{lemma:divisibility}), delivered together with the article's own complete original proof of it (source0.tex lines 528--604, one \texttt{proof} environment of 77 lines). No mathematics is written, repaired or re-derived: statement and proof are the article's own text line for line. Required context retained uncounted: lemma:normal-form-5-CC (lines 472--485); norm-form-CC-2 (lines 475--484). Every definition, display and label of those lines is retained unchanged. Omitted from this package and not redistributed: 1--471, 485--519, 605--1115. Nothing inside a retained excerpt is omitted or altered: every retained body is delivered exactly as the article's lines are written, so no omission receipt of any kind accompanies this package. Citations: the retained excerpts carry no citation command at all, so no bibliography entry of the article is transcribed and no reference is redistributed. Reference adaptation: none was needed and none was made; every reference inside the delivered unit resolves inside it, so no unresolved reference or citation is produced. Printed slips kept literal: no printed word, symbol, hypothesis or number of the counted statement or of its proof was altered. Layout and class: the article's own class is \texttt{amsart} with packages including xcolor, upgreek, amssymb, amsmath, mathtools, amsfonts, amsthm, eucal, fullpage, color, caption, subcaption, enumitem, soul; the wrapper is the lane's pinned \texttt{amsart} editorial wrapper at 10pt on A4, which restates the theorem-like environments the retained text needs and adds the article's own macro definitions that the retained text needs (\texttt{\textbackslash CC}, \texttt{\textbackslash I}, \texttt{\textbackslash J}, \texttt{\textbackslash K}, \texttt{\textbackslash L}, \texttt{\textbackslash PP}, \texttt{\textbackslash SL}, \texttt{\textbackslash cA}, \texttt{\textbackslash cB}, \texttt{\textbackslash overbar}, \texttt{\textbackslash rank}, \texttt{\textbackslash subset}), with extra wrapper packages (bm, bbm, dsfont, xspace, cancel, mathtools, cleveref, enumitem). The delivered numbering is the unit's own. Assets: no retained span contains a figure environment, a graphics-inclusion command, a PDF-page inclusion, a table or a TikZ picture; no image file is copied into this package, the package declares no asset dependency and no image file is redistributed with it. Licence: version arXiv:2507.06942v4 of this article is distributed under the Creative Commons Attribution 4.0 International licence, as recorded in the arXiv licence metadata for this exact version and shown on the abstract page https://arxiv.org/abs/2507.06942v4. A full rights notice is delivered at licenses/arxiv-2507.06942-CC-BY-4.0.txt. Characters: every retained excerpt is pure ASCII, so the delivered engine, XeLaTeX with its default Latin Modern text font, reports no missing character.
\begin{lemma}
    \label{lemma:normal-form-5-CC}
       Let $\overbar{\cA} \in \CC^4 \otimes \bigwedge^2 \CC^5$ be a section such that the scheme $\overbar{C} \subset \PP^3$ cut out by the vanishing of the $4 \times 4$ sub-Pfaffians of $\overbar{\cA}$ has a positive dimensional Zariski tangent space at a point $\overline{p}$. Then there exist unipotent matrices $\overbar{g} \in U_4(\CC) \times U_5(\CC)$, an integer $1 \leq \K \leq 4$, and distinct integers $1 \leq \I,\J,\L \leq 5$ such that if $\overbar \cB = \overbar{g}(\overbar \cA)$ then either:
       \begin{enumerate}[label=(\Alph*)]
           \item \label{norm-form-CC-1} $\overbar{B}_\K = 0$; or
           \item \label{norm-form-CC-2} $\overbar{B}_{k} \neq 0$ for all $k$, and $\bar{b}_{i \I}^{(\K)} = \bar{b}_{i \J}^{(\K)} = \bar{b}_{i \L}^{(\K)} = 0$ for all $1 \leq i \leq 5$, and either:
             \begin{enumerate}[label=(\roman*)]
             \item \label{norm-form-CC-3}
               $\overbar{\mathcal{B}}(\mathbf{x})_{\I\L} = \overbar{\mathcal{B}}(\mathbf{x})_{\J\L} = 0$; or
             \item \label{norm-form-CC-4}
               $\overbar{\mathcal{B}}(\mathbf{x})_{\I\J} = 0$ and the linear forms $\overbar{\mathcal{B}}(\mathbf{x})_{\I\L}$ and $\overbar{\mathcal{B}}(\mathbf{x})_{\J\L}$ are linearly independent.
             \end{enumerate}
        \end{enumerate}
\end{lemma}

\begin{lemma}
  \label{lemma:divisibility}
  Let $R = \CC[u]/u^2$ and let ${\cA} \in R^4 \otimes \bigwedge^2 R^5$ be such that the scheme cut out in $\PP^3_R$ by the $4 \times 4$ sub-Pfaffians of $\cA$ has Zariski tangent space of dimension at least $2$ at a point ${p}$. Then there exists $g \in U_4(\CC) \times U_5(\CC)$ such that either:
  \begin{enumerate}[label=(\Alph*)]
    \item $g(\cA)$ is in normal form of type $(\K)$ at $p$; or 
      \item $g (\cA)$ is in normal form of type $(\I, \J, \K)$ at $p$.
  \end{enumerate}
\end{lemma} 

\begin{proof}
Because $p$ has Zariski tangent space of dimension $\geq 2$, the restriction of the Zariski tangent space to the closed subvariety $u = 0$ has positive dimension and thus we may apply \Cref{lemma:normal-form-5-CC} to $\overbar{\cA} = \cA \pmod{u}$. Suppose that $\cA$ \emph{cannot} be put into normal form of type $(\K)$ so that $\overbar{\cA}$ must be in the form \ref{norm-form-CC-2} from \Cref{lemma:normal-form-5-CC}. In particular we may assume without loss of generality that there exist distinct integers $1 \leq \I,\J, \L \leq 5$ with $\I < \J$ and an integer $1 \leq \K \leq 4$ such that $a_{i \I}^{(\K)} \equiv a_{i \J}^{(\K)} \equiv a_{i \L}^{(\K)} \equiv 0 \pmod u$ for all $1 \leq i \leq 5$ and either:
\begin{enumerate}[label=(\roman*)]
    \item \label{a-mod-u} $\cA(\mathbf{x})_{\I\L} \equiv \cA(\mathbf{x})_{\J\L} \equiv 0 \pmod{u}$; or
    \item \label{b-mod-u} $\cA(\mathbf{x})_{\I\J}  \equiv 0 \pmod{u}$ and $\cA_{I\L}(\mathbf{x})$ and $\cA_{J\L}(\mathbf{x})$ are linearly independent modulo $u$.
\end{enumerate}

In case \ref{a-mod-u}, if $a_{\I\L}^{(\K)} \equiv 0 \pmod{u^2}$ then $\cA$ is already in a normal form of type $(\I, \L, \K)$. Otherwise without loss of generality suppose $\I < \J$ and let $h \in U_5(\CC)$ be the matrix whose $(\I,\J)$-entry is the value of $a_{\J\L}^{(\K)} / a_{\I\L}^{(\K)}$ at $u = 0$ and whose other non-diagonal entries are $0$. After acting by $h$ we may assume that $a_{\J\L}^{(\K)} \equiv 0 \pmod{u^2}$ in which case $\cA$ is in a normal form of type $(\J, \L, \K)$.

In case \ref{b-mod-u} we claim that the equality $a_{\I\J}^{(\K)} \equiv 0 \pmod{u^2}$ holds without any further transformation. To see this, let $\tau = (\tau_4, \tau_5) \in \SL_4(\CC) \times \SL_5(\CC)$ be a pair of permutation matrices for which $\tau_4$ swaps $\K$ and $4$, and $\tau_5$ swaps $\{\I,\J\}$ for $\{4,5\}$. Writing $\mathcal{B} = \tau(\cA)$, we have $\cB(\mathbf{x})_{45} \equiv 0 \pmod u$, and $b_{ij}^{(4)} \equiv 0 \pmod u$ for all $1 \leq i \leq 5$ and $j = 4,5$. That is, $\tau$ swaps the roles of the matrices $A_\K$ and $A_4$ and swaps the roles of the $\{\I,\J\}$ rows with the $\{4,5\}$ rows. In particular we have $a_{\I\J}^{(\K)} = \pm b_{45}^{(4)}$. Under this transformation, the point $p$ is sent to the point ${q} = ([0:0:0:1],0)$ in the scheme in $\PP^3_R$ cut out by the $4 \times 4$ sub-Pfaffians of $\cB$. 

Now consider a transformation $\gamma = (\gamma_4, \gamma_5) \in \SL_4(\CC) \times \SL_5(\CC)$ of the form
\begin{equation*}
    \gamma_4 = \begin{bmatrix}
      * & * & * & * \\
      * & * & * & * \\
      * & * & * & * \\
      0 & 0 & 0 & 1 
    \end{bmatrix}  
    \qquad
    \gamma_5 = \begin{bmatrix}
        * & * & * & 0 & 0 \\
        * & * & * & 0 & 0 \\
        * & * & * & 0 & 0 \\
        0 & 0 & 0 & 1 & 0 \\
        0 & 0 & 0 & 0 & 1 
    \end{bmatrix}.
\end{equation*}
Let $\mathcal{C} = \gamma(\mathcal{B})$ and note that $b_{45}^{(4)} = \pm c_{45}^{(4)}$ remains unchanged up to sign and $([0:0:0:1], 0)$ is a point of the scheme defined by $\mathcal{C}$ whose tangent space has dimension at least $2$. By choosing $\gamma_5$ appropriately, we may assume that modulo $u$ 
\begin{equation*}    
    C_4 \equiv 
    \begin{bmatrix}
      0 & \rho & 0 & 0 & 0 \\
      -\rho & 0 & 0 & 0 & 0 \\
      0 & 0 & 0 & 0 & 0 \\
      0 & 0 & 0 & 0 & 0 \\
      0 & 0 & 0 & 0 & 0 
    \end{bmatrix}
\end{equation*}
and by choosing $\gamma_4$ we may assume that modulo $u$ we have
\begin{equation*}
  C_1 \equiv \begin{bmatrix}
      0 & 0 & * & * & * \\
      0 & 0 & * & * & * \\
      * & * & 0 & \epsilon & 0 \\
      * & * & -\epsilon & 0 & 0 \\
      * & * & 0 & 0 & 0 
  \end{bmatrix}
  \qquad 
  C_2 \equiv \begin{bmatrix}
      0 & 0 & * & * & * \\
      0 & 0 & * & * & * \\
      * & * & 0 & 0 & \lambda \\
      * & * & 0 & 0 & 0 \\
      * & * & -\lambda & 0 & 0 
  \end{bmatrix}
  \qquad
  C_3 \equiv \begin{bmatrix}
      0 & 0 & * & * & * \\
      0 & 0 & * & * & * \\
      * & * & 0 & 0 & 0 \\
      * & * & 0 & 0 & 0 \\
      * & * & 0 & 0 & 0 
  \end{bmatrix}.
\end{equation*}
where $\rho, \epsilon,\lambda \in \{0,1\}$. Because $\rank A_{\K} = 2$ we may assume $\rho = 1$, and because $\cA_{J\L}(\textbf{x})$ and $\cA_{I\L}(\textbf{x})$ are linearly independent we may assume $\epsilon = \lambda = 1$.

Now, a direct calculation shows that on the affine chart where $x_4 \neq 0$ the Jacobian matrix of the $4 \times 4$ sub-Pfaffians of ${\mathcal{C}}(\mathbf{x})$ evaluated at $([0:0:0:1], 0)$ contains the matrix 
\[
\begin{bmatrix}
    0 & \rho \epsilon & 0 \\
    0 & 0 & \rho \lambda \\
    \rho \tfrac{\partial}{\partial u} c_{45}^{(4)} & \rho\tfrac{\partial}{\partial u} c_{34}^{(4)}  & \rho \tfrac{\partial}{\partial u} c_{35}^{(4)}  
\end{bmatrix}
\]
as a full rank submatrix. The scheme defined by $\mathcal{C}$ is assumed to have tangent space of dimension at least $2$ at the point $([0:0:0:1], 0)$ and $\rho = \lambda = \epsilon = 1$. Thus $a_{45}^{(4)} = \pm c_{45}^{(4)} \equiv 0 \pmod{u^2}$, as required.
\end{proof}

\end{document}
